\documentclass[../../main.tex]{subfiles}

\begin{document}

\section{序列的逐点极限与简单逼近}\label{section:ch03-逐点极限简单逼近}

对于具有共同定义域 $E$ 的函数序列 $\{f_{n}\}$ 和 $E$ 上的函数 $f$, 存在几个不同的方式叙述 "序列 $\{f_{n}\}$ 收敛于 $f$", 有必要考虑每一种方式意味着什么.

本章我们考虑具有共同定义域 $E$ 的函数序列 $\{f_{n}\}$ 逐点收敛和一致收敛的概念, 这些概念在初等分析中是熟知的. 在后面的几章我们考虑许多其他模式的函数序列收敛.

\begin{definition}[逐点收敛、a.e. 逐点收敛与一致收敛]\label{definition:ch03-收敛方式}
对具有共同定义域 $E$ 的函数序列 $\{f_{n}\}$, $E$ 上的函数 $f$, 以及 $E$ 的子集 $A$, 我们说
\begin{enumerate}[(i)]
\item 序列 $\{f_{n}\}$ 在 $A$ 上逐点收敛于 $f$, 若
\begin{align*}
\text{对所有 }x\in A,\quad\lim_{n\to\infty}f_{n}(x)=f(x);
\end{align*}
\item 在 $A$ 上序列 $\{f_{n}\}$ a.e. 逐点收敛于 $f$, 若它在 $A\sim B$ 上逐点收敛到 $f$, 其中 $m(B)=0$;
\item 序列 $\{f_{n}\}$ 在 $A$ 上一致收敛于 $f$, 若对每个 $\varepsilon>0$, 存在指标 $N$, 使得
\begin{align*}
\text{对所有 }n\geqslant N,\quad\text{在 }A\text{ 上 }|f-f_{n}|<\varepsilon.
\end{align*}
\end{enumerate}
\end{definition}

当考虑函数序列 $\{f_{n}\}$ 和它们收敛到的函数 $f$ 时, 我们常常隐含假设所有函数有共同的定义域, 我们写 "在 $A$ 上逐点 $\{f_{n}\}\to f$" 以表明在 $A$ 上 $\{f_{n}\}$ 逐点收敛到 $f$, 且对一致收敛用类似的记号.

连续函数的逐点极限不一定是连续的. Riemann 可积函数的逐点极限不一定是 Riemann 可积的. 以下命题首次显示了可测函数有好得多的稳定性.

\begin{proposition}\label{proposition:ch03-逐点极限可测}
令 $\{f_{n}\}$ 为 $E$ 上的 a.e. 逐点收敛于函数 $f$ 的可测函数序列. 则 $f$ 是可测的.
\end{proposition}

\begin{proof}
令 $E_{0}$ 为 $E$ 的子集使得 $m(E_{0})=0$ 而 $\{f_{n}\}$ 在 $E\sim E_{0}$ 上逐点收敛于 $f$. 由于 $m(E_{0})=0$, 从\cref{proposition:ch03-a.e.可测}得知 $f$ 是可测的当且仅当它在 $E\sim E_{0}$ 上的限制是可测的. 因此, 通过可能地用 $E\sim E_{0}$ 替代 $E$, 我们可以假设序列在整个 $E$ 上逐点收敛.

固定数 $c$. 我们必须证明 $\{x\in E\mid f(x)<c\}$ 是可测的. 观察到对于点 $x\in E$, 由于 $\lim_{n\to\infty}f_{n}(x)=f(x)$,
\begin{align*}
f(x)<c
\end{align*}
当且仅当存在自然数 $n$ 和 $k$ 使得对所有 $j\geqslant k$, $f_{j}(x)<c-1/n$.

但是对于任何自然数 $n$ 和 $j$, 集合 $\{x\in E\mid f_{j}(x)<c-1/n\}$ 是可测的, 由于函数 $f_{j}$ 是可测的. 因此, 对任何 $k$, 可测集的可数族的交
\begin{align*}
\bigcap_{j=k}^{\infty}\{x\in E\mid f_{j}(x)<c-1/n\}
\end{align*}
也是可测的. 因此, 由于可测集的可数族的并是可测的,
\begin{align*}
\{x\in E\mid f(x)<c\}=\bigcup_{1\leqslant k,n<\infty}\left[\bigcap_{j=k}^{\infty}\{x\in E\mid f_{j}(x)<c-1/n\}\right]
\end{align*}
是可测的.
\end{proof}

若 $A$ 是任意集合, $A$ 的特征函数 $\chi_{A}$ 是定义为
\begin{align*}
\chi_{A}(x)=
\begin{cases}
1,&\text{若 }x\in A,\\
0,&\text{若 }x\notin A
\end{cases}
\end{align*}
的 $\mathbf{R}$ 上的函数.

显然函数 $\chi_{A}$ 是可测的当且仅当集合 $A$ 是可测的. 因此不可测集的存在性蕴涵着不可测函数的存在性. 可测集的特征函数的线性组合在 Lebesgue 积分中所起的作用类似于阶梯函数在 Riemann 积分中所起的作用, 因而我们将对这些函数命名.

\begin{definition}[简单函数]\label{definition:ch03-简单函数}
定义在可测集 $E$ 上的实值函数 $\varphi$ 称为简单的, 若它是可测的且仅取有限个值.
\end{definition}

我们强调简单函数仅取实值. 简单函数的线性组合与乘积是简单的, 这是由于它们中的每个仅取有限个值. 若 $\varphi$ 是简单的, 具有定义域 $E$ 且取不同值 $c_{1},\cdots,c_{n}$, 则在 $E$ 上
\begin{align*}
\varphi=\sum_{k=1}^{n}c_{k}\cdot\chi_{E_{k}},\quad\text{其中 }E_{k}=\{x\in E\mid\varphi(x)=c_{k}\}.
\end{align*}
$\varphi$ 的特征函数的线性组合的特定的表示称为简单函数 $\varphi$ 的典范表示.

\begin{lemma}[简单逼近引理]\label{lemma:ch03-简单逼近引理}
令 $f$ 为 $E$ 上的可测实值函数. 假设 $f$ 在 $E$ 上有界, 即存在 $M\geqslant 0$, 使得在 $E$ 上 $|f|\leqslant M$. 则对每个 $\varepsilon>0$, 存在定义在 $E$ 上的简单函数 $\varphi_{\varepsilon}$ 和 $\psi_{\varepsilon}$ 具有以下逼近性质:
\begin{align*}
\text{在 }E\text{ 上},\quad\varphi_{\varepsilon}\leqslant f\leqslant\psi_{\varepsilon}\quad\text{且 }0\leqslant\psi_{\varepsilon}-\varphi_{\varepsilon}<\varepsilon.
\end{align*}
\end{lemma}

\begin{proof}
令 $(c,d)$ 为包含 $E$ 的象 $f(E)$ 的开有界区间, 且
\begin{align*}
c=y_{0}<y_{1}<\cdots<y_{n-1}<y_{n}=d
\end{align*}
为闭有界区间 $[c,d]$ 的一个划分, 使得对 $1\leqslant k\leqslant n$, $y_{k}-y_{k-1}<\varepsilon$. 定义
\begin{align*}
I_{k}=[y_{k-1},y_{k})\quad\text{与}\quad E_{k}=f^{-1}(I_{k}),\quad 1\leqslant k\leqslant n.
\end{align*}
由于每个 $I_{k}$ 是区间而函数 $f$ 是可测的, 每个集合 $E_{k}$ 是可测的. 定义 $E$ 上的简单函数 $\varphi_{\varepsilon}$ 和 $\psi_{\varepsilon}$ 为
\begin{align*}
\varphi_{\varepsilon}=\sum_{k=1}^{n}y_{k-1}\cdot\chi_{E_{k}}\quad\text{和}\quad\psi_{\varepsilon}=\sum_{k=1}^{n}y_{k}\cdot\chi_{E_{k}}.
\end{align*}
令 $x$ 属于 $E$. 由于 $f(E)\subseteq(c,d)$, 存在唯一的 $k$ ($1\leqslant k\leqslant n$), 使得 $y_{k-1}\leqslant f(x)<y_{k}$, 因此
\begin{align*}
\varphi_{\varepsilon}(x)=y_{k-1}\leqslant f(x)<y_{k}=\psi_{\varepsilon}(x).
\end{align*}
但 $y_{k}-y_{k-1}<\varepsilon$, 因此 $\varphi_{\varepsilon}$ 和 $\psi_{\varepsilon}$ 有所要求的逼近性质.
\end{proof}

我们添加以下定理到我们已建立的可测函数的刻画中.

\begin{theorem}[简单逼近定理]\label{theorem:ch03-简单逼近定理}
定义在可测集 $E$ 上的扩充实值函数 $f$ 是可测的当且仅当存在 $E$ 上的简单函数序列 $\{\varphi_{n}\}$, 它在 $E$ 上逐点收敛到 $f$ 且具有性质:
\begin{align*}
\text{对所有 }n,\quad\text{在 }E\text{ 上 }|\varphi_{n}|\leqslant|f|.
\end{align*}
若 $f$ 是非负的, 我们可选取 $\{\varphi_{n}\}$ 为递增的.
\end{theorem}

\begin{proof}
由于每个简单函数是可测的, \cref{proposition:ch03-逐点极限可测}告诉我们若一个函数是简单函数序列的逐点极限, 则它是可测的. 剩下来要证明逆命题.

假设 $f$ 是可测的. 我们也假设在 $E$ 上 $f\geqslant 0$. 一般情形通过将 $f$ 表示为非负可测函数的差（见习题 23）得到. 令 $n$ 为自然数. 定义 $E_{n}=\{x\in E\mid f(x)\leqslant n\}$. 则 $E_{n}$ 是可测集而 $f$ 在 $E_{n}$ 的限制是非负有界可测函数. 将简单逼近引理用于 $f$ 在 $E_{n}$ 的限制, 且选取 $\varepsilon=1/n$, 我们可以选取定义在 $E_{n}$ 上的简单函数 $\varphi_{n}$ 和 $\psi_{n}$, 它们具有以下逼近性质:
\begin{align*}
\text{在 }E_{n}\text{ 上},\quad 0\leqslant\varphi_{n}\leqslant f\leqslant\psi_{n}\text{ 且 }0\leqslant\psi_{n}-\varphi_{n}<1/n.
\end{align*}
观察到
\begin{align}
\text{在 }E_{n}\text{ 上},\quad 0\leqslant\varphi_{n}\leqslant f\text{ 且 }0\leqslant f-\varphi_{n}\leqslant\psi_{n}-\varphi_{n}<1/n.\label{align:ch03-逼近估计}
\end{align}
若 $f(x)>n$, 通过设 $\varphi_{n}(x)=n$, 将 $\varphi_{n}$ 延拓到整个 $E$. 函数 $\varphi_{n}$ 是定义在 $E$ 上的简单函数且在 $E$ 上 $0\leqslant\varphi_{n}\leqslant f$. 我们宣称序列 $\{\varphi_{n}\}$ 在 $E$ 上逐点收敛于 $f$. 令 $x$ 属于 $E$.

\textbf{情形 1:} 假设 $f(x)$ 是有限的. 选取一个自然数 $N$ 使得 $f(x)<N$. 则对于 $n\geqslant N$,
\begin{align*}
0\leqslant f(x)-\varphi_{n}(x)<1/n.
\end{align*}
因此 $\lim_{n\to\infty}\varphi_{n}(x)=f(x)$.

\textbf{情形 2:} 假设 $f(x)=\infty$. 则对所有 $n$, $\varphi_{n}(x)=n$, 因此 $\lim_{n\to\infty}\varphi_{n}(x)=f(x)$.

通过用 $\max\{\varphi_{1},\cdots,\varphi_{n}\}$ 代替每个 $\varphi_{n}$, 我们有 $\{\varphi_{n}\}$ 是递增的.
\end{proof}

\begin{exercise}\label{exercise:ch03-12}
令 $f$ 为 $E$ 上的有界可测函数. 证明存在 $E$ 上的简单函数序列 $\{\varphi_{n}\}$ 和 $\{\psi_{n}\}$, 使得 $\{\varphi_{n}\}$ 是递增的而 $\{\psi_{n}\}$ 是递减的, 且这些序列的每个在 $E$ 上一致收敛于 $f$.
\end{exercise}

\begin{exercise}\label{exercise:ch03-13}
实值可测函数称为半简单的, 若它仅取可数个值. 令 $f$ 为 $E$ 上的任意可测函数. 证明存在 $E$ 上的半简单函数序列 $\{f_{n}\}$, 它在 $E$ 上一致收敛于 $f$.
\end{exercise}

\begin{exercise}\label{exercise:ch03-14}
令 $f$ 为 $E$ 上的可测函数, 其在 $E$ 上 a.e. 有限且 $m(E)<\infty$. 证明对每个 $\varepsilon>0$, 存在包含于 $E$ 的可测集 $F$, 使得 $f$ 在 $F$ 上是有界的且 $m(E\sim F)<\varepsilon$.
\end{exercise}

\begin{exercise}\label{exercise:ch03-15}
令 $f$ 为 $E$ 上的有界可测函数, 其在 $E$ 上 a.e. 有限且 $m(E)<\infty$. 证明对每个 $\varepsilon>0$, 存在包含于 $E$ 的可测集 $F$ 和 $E$ 上的简单函数序列 $\{\varphi_{n}\}$ 使得在 $F$ 上一致地 $\{\varphi_{n}\}\to f$ 且 $m(E\sim F)<\varepsilon$.（提示: 见前一个习题.）
\end{exercise}

\begin{exercise}\label{exercise:ch03-16}
令 $I$ 为闭有界区间而 $E$ 是 $I$ 的可测子集. 令 $\varepsilon>0$. 证明存在 $I$ 上的阶梯函数 $h$ 和 $I$ 的可测子集 $F$, 使得
\begin{align*}
\text{在 }F\text{ 上},\quad h=\chi_{E}\text{ 且 }m(I\sim F)<\varepsilon.
\end{align*}
（提示: 用第 2 章的\cref{theorem:ch02-有限区间逼近}.）
\end{exercise}

\begin{exercise}\label{exercise:ch03-17}
令 $I$ 为闭有界区间而 $\psi$ 是定义在 $I$ 上的简单函数. 令 $\varepsilon>0$. 证明存在 $I$ 上的阶梯函数 $h$ 和 $I$ 的可测子集 $F$, 使得
\begin{align*}
\text{在 }F\text{ 上},\quad h=\psi\text{ 且 }m(I\sim F)<\varepsilon.
\end{align*}
（提示: 用简单函数是特征函数的线性组合的事实以及前一个习题.）
\end{exercise}

\begin{exercise}\label{exercise:ch03-18}
令 $I$ 为闭有界区间而 $f$ 是定义在 $I$ 上的有界可测函数. 令 $\varepsilon>0$. 证明存在 $I$ 上的阶梯函数 $h$ 和 $I$ 的可测子集 $F$, 使得
\begin{align*}
\text{在 }F\text{ 上},\quad|h-f|<\varepsilon\text{ 且 }m(I\sim F)<\varepsilon.
\end{align*}
\end{exercise}

\begin{exercise}\label{exercise:ch03-19}
证明如同 max 与 min, 两个简单函数的和与积是简单的.
\end{exercise}

\begin{exercise}\label{exercise:ch03-20}
令 $A$ 和 $B$ 为任意集. 证明
\begin{align*}
\chi_{A\cap B}=\chi_{A}\cdot\chi_{B},\quad\chi_{A\cup B}=\chi_{A}+\chi_{B}-\chi_{A}\cdot\chi_{B},\quad\chi_{A^{c}}=1-\chi_{A}.
\end{align*}
\end{exercise}

\begin{exercise}\label{exercise:ch03-21}
对于具有共同定义域的可测函数序列 $\{f_{n}\}$, 证明以下函数都是可测的:
\begin{align*}
\inf\{f_{n}\},\quad\sup\{f_{n}\},\quad\liminf\{f_{n}\}\text{ 与 }\limsup\{f_{n}\}.
\end{align*}
\end{exercise}

\begin{exercise}\label{exercise:ch03-22}
（Dini 定理）令 $\{f_{n}\}$ 为 $[a,b]$ 上的连续函数的递增序列, 它在 $[a,b]$ 上逐点收敛于 $[a,b]$ 上的连续函数 $f$. 证明该收敛在 $[a,b]$ 上是一致的.（提示: 令 $\varepsilon>0$. 对每个自然数 $n$, 定义 $E_{n}=\{x\in[a,b]\mid f(x)-f_{n}(x)<\varepsilon\}$. 证明 $\{E_{n}\}$ 是 $[a,b]$ 的开覆盖并用 Heine-Borel 定理.）
\end{exercise}

\begin{exercise}\label{exercise:ch03-23}
将可测函数表示为非负可测函数的差, 进而基于非负可测函数的特殊情形证明一般的简单逼近定理.
\end{exercise}

\begin{exercise}\label{exercise:ch03-24}
令 $I$ 为区间而 $f:I\to\mathbf{R}$ 是递增的. 通过首先证明对每个自然数 $n$ 严格递增的函数 $x\mapsto f(x)+x/n$ 是可测的, 接着取逐点极限, 证明 $f$ 是可测的.
\end{exercise}

\end{document}